% Expanded decision chapter: the shared decision and reader journey
\chapter{What the first investment would buy}
\label{sec:opening}

The case and the wider record now let us state the investment plainly. A small
team would have twelve weeks to build a review system, try it on one
consequential document, and measure whether it saves the attention of the
person who has to decide what the document means. The people making that
decision can assess a research design and a financial case; they should not
have to reconstruct the project's private history before they can judge it.

The practical question is straightforward: is this bounded effort worth
funding, given what the cases and the literature already tell us and what the
first build still has to learn?

\section{The first investment}

\hv{} is a reader-first revision system for technical documents. It does not
try to make prose sound generically ``human,'' and it does not replace the
author's judgment with a detector or a model score. It locates problems that
can be shown in the source, protects equations, numbers, citations, and claims
while a draft changes, and ends with a decision from a named human reader.

The first investment is limited. Fund a twelve-week MVP, a small evaluation
corpus, and one live-document feasibility run. At the end of that work, we
should decide whether a comparative study against the current
combination of linters, general language-model editing, version control, and
manual review is worth funding. This request does not fund general
deployment or support a claim that \hv{} already improves expert writing.

The eventual success criterion is economic as well as technical: fewer
human-feedback rounds to reader acceptance, without increasing reader time or
changes to protected meaning. If the system makes review slower, creates too
many low-value alerts, or cannot show what happened to a substantive change,
we should stop or redesign it even if its surface scores improve.

\begin{table}[H]
\centering
\small
\begin{tabular}{@{}p{0.25\textwidth}p{0.68\textwidth}@{}}
\toprule
Decision question & Proposed answer \\
\midrule
What problem is worth paying to solve? & Repeated diagnosis, revision, and rereading when a document is correct but leaves an expert reader to reconstruct its purpose and argument. \\
What would we build first? & A local-first parser, located diagnostics, protected comparison, reader review, and an evaluation corpus for LaTeX and Markdown. \\
What would count as evidence? & A working end-to-end case followed, if feasible, by a preregistered comparison of feedback rounds, reader time, and meaning preservation. \\
What is not yet established? & That this workflow improves expert economics or research writing, or that any detector can decide whether a document is acceptable. \\
What decision is requested now? & Fund the bounded MVP and feasibility case; defer deployment and efficacy claims until the comparison earns them. \\
\bottomrule
\end{tabular}
\caption{The investment decision in one page. The later chapters supply the evidence and design behind each answer.}
\label{tab:opening-decision}
\end{table}

\Needspace{4.5cm}
\begin{decisionbox}{The recommendation}
Fund the bounded build and the feasibility case. Do not treat deployment or a
passing surface score as evidence that expert writing has improved. Fund a
larger study only if the workflow is usable, protects meaning, and leaves the
reader outcome credible.
\end{decisionbox}

\section{What we know, and what the first build must learn}

The work is not starting from a blank page. The internal cases provide a
concrete failure and repair vocabulary; the literature supplies mechanisms and
limits; and the software review identifies components that can be tested rather
than invented. What remains unknown is the value of connecting those pieces for
technical authors and their readers.

\begin{table}[H]
\centering
\small
\begin{tabularx}{\textwidth}{@{}p{0.24\textwidth}L@{}}
\toprule
\textbf{Already established} & \textbf{What the record supports} \\
\midrule
The costly failure & A document can compile and pass agent review while an expert unfamiliar with the project cannot state its purpose or requested action. \\
The mechanism & Internal vocabulary, mechanical repetition, defensive prose, and missing substance mask one another in a predictable order. \\
The available pieces & Parsers, linters, protected comparison, citation tools, and reader-review methods exist, but no validated system connects them. \\
\midrule
\textbf{Still missing} & \textbf{What the MVP must learn} \\
\midrule
Economic value & Whether the connected workflow reduces feedback and rereading time enough to justify its cost. \\
Generalization & Whether the result survives new writers, genres, and readers rather than one rescued document. \\
Efficacy & Whether a comparative study can show fewer feedback rounds without meaning violations or slower readers. \\
\bottomrule
\end{tabularx}
\caption{The current evidence separates feasibility from efficacy.}
\label{tab:opening-status}
\end{table}

\section{What the first system must get right}

The alternatives already in use each solve one part of the problem. A grammar
checker catches local errors; a general language model drafts and rewrites; a
detector estimates whether text resembles a model; and a version-control
system records that a file changed. An investable product has to satisfy five
conditions at once, because improving one layer while damaging another merely
moves the reader's cost.

\begin{table}[H]
\centering
\small
\begin{tabularx}{\textwidth}{@{}p{0.23\textwidth}p{0.31\textwidth}L@{}}
\toprule
Condition & Why it matters & First-release test \\
\midrule
\textbf{Reader outcome} & The document exists to support a decision, not to
  win a style score. & A named reader recovers the problem, evidence,
  qualification, and requested action from the rendered document. \\
\textbf{Meaning protection} & A fluent edit is a failure if it changes an
  equation, number, citation, assumption, or conclusion. & A protected
  comparison explains every change to those objects and blocks silent loss. \\
\textbf{Located diagnosis} & Authors can evaluate a finding only when they
  can see where it came from and why it matters. & Each alert names a span,
  category, evidence, and editorial question; no alert edits the source. \\
\textbf{Local provenance} & Technical manuscripts may contain confidential
  research and unpublished decisions. & The local-first path reproduces a
  finding from a source snapshot without sending the manuscript to a remote
  service by default. \\
\textbf{Economic discipline} & A tool that adds review work is not an
  improvement, however polished its output. & The feasibility run measures
  feedback rounds, reader time, false-positive burden, and meaning errors
  before any scale-up decision. \\
\bottomrule
\end{tabularx}
\caption{Five conditions for a useful first release.}
\label{tab:opening-conditions}
\end{table}

The comparison is therefore with a bundle of existing practices, not with a
single rival. Humanvoice earns a place only if the bundle becomes easier to
inspect and cheaper to use without weakening the author's control. The later
chapters test each condition against research findings, source-level package
reviews, and the project's own failure record.

\section{The cost of a document that makes the reader do the work}

A technical document can compile without warnings, survive several model
reviews, and contain correct mathematics while failing in its first few
minutes with a human reader. One funding proposal in this project's record
did exactly that. Its reader could not state the purpose, could not tell why
private file names and phase labels appeared in the argument, and could not
identify the action being requested. The failure was not a missing comma. The
document had transferred the work of constructing the argument from its author
to a scarce expert reader. The preceding worked case makes that transfer
concrete; this section draws out its economic consequence rather than treating
the case as an external efficacy study.

That transfer has a simple economic consequence. A rejected draft starts
another cycle of diagnosis, revision, and rereading. The author spends time,
but so does the person whose judgment the document was meant to support. For a
research group, investment committee, or executive team, the scarce input is
often not another paragraph of writing. It is the uninterrupted attention of
someone who can decide whether the work deserves resources or a change in
direction.

The current toolchain lowers parts of this cost. Grammar checkers find local
errors. Prose linters expose repeated phrases. A general language model can
draft, rewrite, or comment. A citation service can confirm that a DOI exists,
and version control can show that text changed. None of these, on its own,
answers the decision question: can the intended reader now reconstruct the
problem, the mechanism, the evidence, and the requested action? A fluent
rewrite can also change a number, assumption, citation, or conclusion while
making the sentence look better.

\subsection{A concrete repair}

The product's promise is easier to see in a small, recorded change. In one
internal economics manuscript, a sentence carried a project-status word that
made sense to its authors but not to an unaided reader:

\begin{repairpair}
\before{``It translates one declared restricted reconstruction into Dynare and
compares like with like, including likelihood components, observations,
counterfactual paths, and lower-bound conditions.''}
\after{``It translates the restricted reconstruction into Dynare and makes a
like-for-like comparison of likelihood components, observations,
counterfactual paths, and lower-bound conditions.''}
\end{repairpair}

The repair removes the private status marker and makes the comparison the
sentence's action. It does not change an equation, a result, or a qualification.
An MVP should make this kind of located, reviewable change easy to propose and
safe to reject; a generic score would not show either fact.

\section{What the product does in one document}

\hv{} sits beside the author's existing editor and repository. A document
moves through a short, observable sequence.

First, the author identifies the intended reader and the decision the document
must support. The system parses a LaTeX or Markdown source into ordinary prose
and protected content. Equations, numbers, labels, citations, quotations, and
declared claims retain their source locations.

Second, it presents a small ranked set of findings. Some are deterministic: a
private path has entered the manuscript, several sections open in the same
way, a citation's metadata does not match, or a protected number changed.
Others are editorial questions attached to a span: does this qualification
prevent a real misunderstanding, or only protect the author? Is the mechanism
that connects the paragraph's evidence to its conclusion missing? Findings
remain findings. The system does not silently edit the source.

Third, the author or an explicitly chosen editor revises the draft. A
protected comparison reports what happened to each equation, number, label,
citation, quotation, and declared claim. A substantive change is allowed, but
it cannot disappear inside a smoother sentence.

Finally, a reader who has not seen the project history receives the rendered
document without the context that made its private language seem normal. The reader
answers a short set of questions tied to the document's purpose: What problem
is being solved? What is proposed? Why is the obvious alternative
insufficient? What evidence matters? What action is requested? A named human
remains the final authority. A missing answer sends the document back for
substantive revision, not another round of synonym replacement.

\begin{quote}\small\ttfamily
source -> parse and protect -> located findings\\
-> author revision -> protected comparison\\
-> reader without project history -> human decision
\end{quote}

\begin{figure}[H]
\centering
\begin{tikzpicture}[node distance=2mm and 2mm,font=\small,scale=0.86,transform shape]
  \node[hvstage] (source) {Source document};
  \node[hvstage,right=of source] (parse) {Parse source\\and protect meaning};
  \node[hvstage,right=of parse] (find) {Located findings\\with evidence};
  \node[hvstage,right=of find] (revise) {Author revision\\and protected comparison};
  \node[hvstage,right=of revise] (read) {Independent read\\and human decision};
  \draw[hvflow] (source) -- (parse);
  \draw[hvflow] (parse) -- (find);
  \draw[hvflow] (find) -- (revise);
  \draw[hvflow] (revise) -- (read);
\end{tikzpicture}
\caption{The product connects measurement to judgment without making either one impersonate the other.}
\label{fig:opening-journey}
\end{figure}

The reader sees one manuscript. The project retains only the provenance needed
to reproduce a finding, understand a revision, and evaluate the workflow. That
separation is an implementation boundary, not a second reader-facing
document: this volume carries the argument, while machine records remain in
the repository for reproducibility.

\section{What the evidence allows us to promise}

The research literature makes the design plausible and limits its promise.
Controlled experiments find that generative assistance can make some writing
and knowledge-work tasks faster or better scored \citep{noy2023experimental}.
Field studies report productivity changes in workplace settings, but average
effects vary by task and worker experience \citep{dillon2025shifting,brynjolfsson2025work}.
Assistance can even reduce correctness when a task lies outside a model's
capability frontier \citep{dellacqua2026jagged}. These studies establish that
assistance changes work; they do not estimate the value of this workflow for
expert research writing.

Studies of feedback evaluate comments by the revisions they induce, not by
their fluency. Models can identify local defects while missing the central
problem in a passage \citep{nair2024closing,rashkin2025story}. Systematic
reviews find more consistent gains in grammar, process, coherence, and surface
accuracy than in argumentation, discourse organization, or creativity
\citep{meng2026systematic,abudalfa2026awe}. Model judges are useful screening
instruments, but their preferences change with order, length, model family,
and self-preference \citep{zheng2023judging,panickssery2024llm,dubois2024length}.
Disclosure can change reader response, assistance can increase individual
performance while reducing collective diversity, and a DOI match does not
show that a paper supports the sentence citing it
\citep{li2024disclosure,doshi2024generative,padmakumar2024writing,mugaanyi2024citation}.

Together these findings support a narrow product hypothesis: connecting
located diagnosis, protected editing, and direct reader evidence may reduce
wasted revision and rereading. They do not justify autonomous rewriting, an
authorship detector, a scalar ``human-likeness'' score, or an efficacy claim
before the pilot.

\section{What the product joins together}

The proposition is an interface claim, not a claim that one new model will
write better sentences. Humanvoice should join four activities that are
usually separated: locating a problem in the source, preserving the objects
that carry meaning, asking the author to choose a repair, and measuring what a
reader can recover afterward.

\begin{enumerate}[label=\textbf{C\arabic*.},leftmargin=*]
\item \textbf{A bounded diagnosis.} Findings have a location, a category, and
  an explanation. They are useful because an author can disagree with them.
\item \textbf{A protected revision.} Equations, numbers, citations, labels,
  quotations, and claims remain visible while prose changes.
\item \textbf{A reader outcome.} The endpoint is whether a named reader can
  reconstruct the problem, mechanism, evidence, and requested action.
\item \textbf{An economic test.} The workflow earns further investment only
  when it reduces costly feedback without increasing reader time or meaning
  violations.
\end{enumerate}

\begin{takeawaybox}{Why this is a product rather than another writing guide}
The first three claims describe a usable workflow. The fourth makes it
falsifiable. A policy can recommend these behaviors; a product can record them,
replay them, and discover whether they save expert attention.
\end{takeawaybox}

\section{The economic test}

The value of the project should be calculated from its own documents and
readers. Let $r_0$ and $r_1$ be the number of feedback rounds under the current
and proposed workflows, $t_0$ and $t_1$ the expert-reader hours per round,
$c_r$ the loaded cost of that time, $N$ the annual number of relevant
documents, and $C$ the annualized build and maintenance cost. The directly
measurable reader-time component of annual value is
\begin{equation}
V = N(r_0t_0-r_1t_1)c_r-C.
\label{eq:opening-value}
\end{equation}

The equation is a decision aid, not a forecast. The project does not yet have
defensible values for all of its inputs, and no published study supplies them.
The feasibility case should measure the baseline distribution of rounds,
reader time, false-positive burden, and meaning-preservation failures. A later
comparison can then price the investment honestly and report results by writer
experience, document genre, and reader role.

\section{What the first release enables}

The same workflow can serve several high-value document settings without
pretending that they are the same genre:

\begin{description}
\item[Research proposals.] Keep the mechanism and requested decision visible
  while evidence and caveats are revised.
\item[Technical monographs and model notes.] Protect notation, equations, and
  calibration claims while a long argument is made easier to teach.
\item[Executive decision briefs.] Compress the route to an action without
  turning the technical appendix into the main story.
\item[Implementation and evaluation records.] Preserve enough provenance to
  reproduce a finding without asking the decision reader to read the ledger.
\end{description}

These are application targets, not evidence that the workflow already works in
each setting. The feasibility case should choose one live document and one
named reader first; generalization is a result to earn.

\section{Where the evidence leads}
\label{sec:route}

The next chapters examine the research that explains machine-typical language,
assistance, feedback, judging, trust, and citation reliability. They then
examine software at source level where possible and distinguish operational
suitability from effectiveness. The argument returns to the policies and tools
already in this project before turning those findings into a product boundary,
a feasible implementation, and a testable evaluation.

The internal case record appears once because it is useful evidence about
failure mechanisms and test fixtures, but it is not a substitute for external
research. The research chapters identify mechanisms and limits; the software
chapters ask which components can perform a named job; and the product chapters
turn those answers into a protected revision path. The final chapters ask
whether the path saves attention and what result would justify the next
investment. A result that fails those tests tells us to stop or to narrow the
claim.
